Large language models (LLMs) now mediate a growing fraction of consequential decisions in healthcare, law, education, scientific writing, and government. The frontier systems span GPT-3 (Brown et al., 2020) at 175B parameters, GPT-4 (OpenAI, 2023), and PaLM 2 (Anil et al., 2023), with open-weight families LLaMA-2 (Touvron et al., 2023) at 7B/13B/70B and LLaMA-3.1 (Meta, 2024) at 8B/70B/405B. Claude 3 Sonnet (Anthropic, 2024), Mistral 7B (Jiang et al., 2023), and Gemma-2 (Google, 2024) at 2B/9B/27B round out the lineup. Despite this capability explosion, these models remain notoriously opaque. A Llama-3.1 70B instance computes its next-token distribution by composing roughly 80 transformer layers of attention and MLP operations across an 8,192-dimensional residual stream. Its causal antecedents are unavailable to direct inspection. This opacity is not a cosmetic inconvenience. It blocks medical deployment under EU AI Act traceability rules. It hampers safety auditing of jailbreak vulnerabilities. It undermines scientific reproducibility when the model is used as a measurement instrument. It prevents users from calibrating trust on individual outputs. Explainability for large language models renders an LLM's predictions, internal computations, or generated reasoning intelligible to human stakeholders. It is therefore one of the most active research frontiers in the post-2020 NLP landscape. This survey synthesises the rapid progression from gradient-based saliency on early Tr\ldots{}
